\section{Análise de Requisitos}

O presente projeto visa o desenvolvimento de um sistema para monitoramento de temperatura e umidade, utilizando o microcontrolador ESP32. O principal objetivo é garantir o controle da temperatura e umidade por meio da coleta de dados em tempo real com um sensor dedicado.

O sistema foi concebido para fornecer tanto interação local quanto remota. Localmente, o dispositivo informa o usuário por meio de uma tela física e indicadores visuais (LEDs). Remotamente, tira proveito da conectividade Wi-Fi do ESP32 para disponibilizar uma interface web dinâmica. Por meio desse painel de controle web, o usuário pode visualizar a telemetria e gerenciar o comportamento do hardware, alternando entre diferentes modos de operação.

A listagem a seguir documenta as necessidades do sistema, divididas em Requisitos Funcionais (que descrevem as funcionalidades e ações que o sistema deve realizar) e Requisitos Não-Funcionais (que definem as restrições, o desempenho e os critérios de qualidade esperados para a arquitetura).

\subsection{Requisitos Funcionais}

\subsubsection{[RF01] Medição de Temperatura}
O sistema deve ser capaz de coletar os dados de temperatura do ambiente em tempo real através de um sensor. \\
\textbf{Prioridade:} Essencial

\subsubsection{[RF02] Medição de Umidade}
O sistema deve ser capaz de coletar os dados de umidade relativa do ar em tempo real através de um sensor dedicado. \\
\textbf{Prioridade:} Essencial

\subsubsection{[RF03] Exibição em Tempo Real via Web}
O sistema deve disponibilizar as leituras de temperatura e umidade em uma interface web com atualização contínua. \\
\textbf{Prioridade:} Essencial

\subsubsection{[RF04] Alternância entre Modos de Operação}
O sistema deve permitir que o usuário alterne entre três modos de operação pela interface web. \\
\textbf{Prioridade:} Essencial

\subsubsection{[RF05] Modo Desligado}
O sistema deve desativar a coleta do sensor e todos os dispositivos de saída, quando o modo Desligado estiver selecionado, informando o estado no display e na interface web. \\
\textbf{Prioridade:} Essencial

\subsubsection{[RF06] Modo Ligado}
O sistema deve manter a leitura do sensor ativa, enviando os dados em tempo real para a interface web, e acionar o buzzer e o ventilador somente mediante autorização do usuário. \\
\textbf{Prioridade:} Essencial

\subsubsection{[RF07] Modo Automático}
O sistema deve realizar a leitura autônoma dos dados sem precisar de uma ação manual do usuário. \\
\textbf{Prioridade:} Essencial

\subsubsection{[RF08] Indicação Visual por LEDs no Hardware}
O circuito físico deve conter LEDs que reflitam o modo de operação ativo. \\
\textbf{Prioridade:} Essencial

\subsubsection{[RF09] Alertas e Notificações}
O sistema deve alertar visualmente na interface e disparar notificação quando temperatura ou umidade ultrapassarem limites configurados. \\
\textbf{Prioridade:} Desejável

\subsubsection{[RF10] Comunicação sem Fio (Wi-Fi)}
A integração entre o ESP32 e a interface web deve ocorrer por meio de rede local sem fio. \\
\textbf{Prioridade:} Essencial

\subsubsection{[RF11] Exibição em Display}
O sistema deve conter um display físico no hardware capaz de exibir localmente a temperatura e a umidade medidas em tempo real. \\
\textbf{Prioridade:} Essencial

\subsubsection{[RF12] Coletar dados}
O sistema deve coletar os dados de temperatura através de um sensor. \\
\textbf{Prioridade:} Essencial

\subsubsection{[RF13] Exibir informações}
O sistema deve apresentar a temperatura no display. \\
\textbf{Prioridade:} Essencial

\subsubsection{[RF14] Acionamento do Ventilador}
O sistema deve acionar o ventilador a partir de uma temperatura elevada. \\
\textbf{Prioridade:} Desejável

\subsection{Requisitos Não-Funcionais}

\subsubsection{[NF01] Tempo de Resposta da Interface}
O sistema deve garantir um tempo de resposta inferior a 2 segundos para a atualização dos dados na interface web após a coleta. \\
\textbf{Prioridade:} Essencial

\subsubsection{[NF02] Precisão na medição de temperatura}
O sensor utilizado deve possuir uma margem de erro baixa para garantir a confiabilidade no monitoramento da temperatura do ambiente. \\
\textbf{Prioridade:} Essencial

\subsubsection{[NF03] Precisão na medição de umidade}
O sensor utilizado deve apresentar margem de erro aceitável para garantir a fidelidade dos dados de umidade. \\
\textbf{Prioridade:} Essencial

\subsubsection{[NF04] Ter precisão}
O sistema realizará as medições de temperatura adequada ao objetivo do projeto. \\
\textbf{Prioridade:} Essencial

\subsubsection{[NF05] Consumir pouca energia}
O sistema deve consumir pouca energia. \\
\textbf{Prioridade:} Essencial